% !TEX root = ../main.tex
\section{SQL Templates, Configuration, and Extended Notes}
\label{ch:appendix-sql}

This appendix collects representative SQL patterns, configuration keys, and extended methodological notes that support reproduction.

\dissertationSubheading[sec:sql-approval]{Approval log extraction pattern}

ERC-20 \texttt{Approval} events are identified by topic~$0$ equal to the keccak256 hash of \texttt{Approval(address,address,uint256)}. A simplified monthly query pattern is:

\begin{verbatim}
SELECT
  block_timestamp,
  block_number,
  transaction_hash,
  log_index,
  address AS token,
  topics[SAFE_OFFSET(1)] AS owner_topic,
  topics[SAFE_OFFSET(2)] AS spender_topic,
  data AS value_data
FROM `bigquery-public-data.goog_blockchain_arbitrum_one_us.logs`
WHERE block_timestamp >= TIMESTAMP('2025-12-01')
  AND block_timestamp <  TIMESTAMP('2026-01-01')
  AND topics[SAFE_OFFSET(0)] = '0x8c5be1e5ebec7d5bd14f71427d1e84f3dd0314c0f7b2291e5b200ac8c7c3b925'
  AND (
    LOWER(CONCAT('0x', SUBSTR(topics[SAFE_OFFSET(1)], 27))) IN UNNEST(@wallet_list)
    OR LOWER(CONCAT('0x', SUBSTR(topics[SAFE_OFFSET(2)], 27))) IN UNNEST(@wallet_list)
  );
\end{verbatim}

Wallet filtering via \texttt{@wallet\_list} keeps scans proportional to the cohort rather than the full chain. Monthly batching respects per-query byte caps documented in the study's acquisition record.

\dissertationSubheading[sec:sql-transfer]{Transfer log extraction pattern}

\texttt{Transfer} events use topic~$0$ for \texttt{Transfer(address,address,uint256)}:

\begin{verbatim}
SELECT
  block_timestamp,
  block_number,
  transaction_hash,
  log_index,
  address AS token,
  topics[SAFE_OFFSET(1)] AS from_topic,
  topics[SAFE_OFFSET(2)] AS to_topic,
  data AS value_data
FROM `bigquery-public-data.goog_blockchain_arbitrum_one_us.logs`
WHERE block_timestamp >= TIMESTAMP('2025-12-01')
  AND block_timestamp <  TIMESTAMP('2026-01-01')
  AND topics[SAFE_OFFSET(0)] = '0xddf252ad1be2c89b69c2b068fc378daa952ba7f163c4a11628f55a4df523b3ef'
  AND (
    LOWER(CONCAT('0x', SUBSTR(topics[SAFE_OFFSET(1)], 27))) IN UNNEST(@wallet_list)
    OR LOWER(CONCAT('0x', SUBSTR(topics[SAFE_OFFSET(2)], 27))) IN UNNEST(@wallet_list)
  );
\end{verbatim}

Post-processing applies token decimals, builds time-decayed edge weights, and restricts nodes to the evaluation sample before PageRank.

\dissertationSubheading[sec:config-keys]{Configuration keys}

The pipeline's parameters are collected in a single versioned configuration file. Salient keys include:

\begin{itemize}
\item observation window start and end timestamps;
\item PageRank damping, tolerance, and maximum iterations;
\item logistic decay $k$ and $t_0$;
\item benchmark repetition count and subsample seeds;
\item output locations for intermediate artifacts and table fragments (\texttt{results/tables/}).
\end{itemize}

If the damping factor or the decay parameter are changed, then the entire evaluation has to be run again (including robustness sweeps and exporting the tables) and all figures re-created from the same artifacts so that the numerics match what is reported in the pipeline's machine-readable summary.

\dissertationSubheading[sec:complexity-notes]{Complexity notes}

Let $n$ be the number of nodes and $m=|E|$ the number of edges. Constructing the latest-allowance edges takes $O(R_a\log R_a)$ time if the $R_a$ approval rows are already sorted by timestamp, or $O(R_a)$ expected time using a hash map keyed on (token, owner, spender). Constructing the transfer edges takes $O(R_t)$ time over the transfer rows. PageRank runs in $O(m + mI)$ time for $I$ iterations, where the first term is the cost of building the sparse operator. The iteration count and edge counts for both methods are shown in Table~\ref{tab:benchmark-runtime}. EndorseRank uses about one-third the number of iterations that AWP does and has $m_{\mathrm{approve}}\ll m_{\mathrm{transfer}}$, which means that both aspects work in EndorseRank's favor.

Memory usage is dominated by the edge lists and score vectors. The peak-memory ratio in Table~\ref{tab:benchmark-runtime} is on the same order of magnitude as the edge-count ratio.

\dissertationSubheading[sec:interpretation-checklist]{Conventions for the result tables}

The tables of Chapter~\ref{ch:results} follow these conventions:

\begin{enumerate}
\item Indicate whether a row corresponds to a social method (EndorseRank or AWP).
\item Associate each family with its target concept (allowance vs.\ transfer).
\item Favor relative comparisons on the same cohort rather than absolute $\tau$ thresholds.
\item Treat runtime tables as representing PageRank solve time on pre-built edge lists, not end-to-end ETL time, and assume that the numbers will be the same whenever the same graph is reported in the same configuration (see Section~\ref{sec:computational-benchmark}).
\item Report every $\tau$ with its interval and derive differences between coefficients from Table~\ref{tab:tau-diff} rather than subtracting the point estimates.
\item Treat the sample-definition sweep (Table~\ref{tab:robustness-sample-size}) as a sensitivity analysis and do not compare its absolute values to the matched-cohort tables.
\end{enumerate}

\dissertationSubheading[sec:archived-baselines]{Archived extended baselines}

The three-method (including the now-withdrawn GF-PR reference), seven-method and six-Aave alignment matrices have been archived with the evaluation pipeline, along with the corresponding CSV matrices. They are excluded from the main comparison because some variants rely on additional event types (such as credit delegation) or introduce concepts that are not part of the EndorseRank--AWP question.

\dissertationSubheading[sec:limitations-reprise]{Limitations}

The SQL templates do not mean that other chains were queried, nor do the archived baselines serve as peer comparisons in this study. The numerics presented are those in the pipeline's machine-readable summary for the matched cohort ($n=5{,}521$) and the expanded wallet pool ($N=99{,}080$).
